% SEGMENT OLP-0636-S01
% Part: history
% Chapter: set-theory
% Section: pathologies

\documentclass[../../../include/open-logic-section]{subfiles}

\begin{document}

\olfileid{his}{set}{pathology}
\olsection{வழக்கத்திற்கு மாறான எடுத்துக்காட்டுகள்}

ஆனால் \emph{எல்லை} என்ற வரையறை, வழக்கத்திற்கு மாறான சில
கட்டுமானங்களையும் அனுமதிப்பதாகத் தெரியவந்தது.

1830-களில் போல்சானோ, \emph{எல்லா இடங்களிலும் தொடர்ச்சியாக} இருந்து,
\emph{எந்த இடத்திலும் வகையிடத்தக்கதாக இல்லாத} ஒரு சார்பைக் கண்டறிந்தார்.
(வருந்தத்தக்க வகையில், போல்சானோ இதை வெளியிடவில்லை; வெயர்ஸ்ட்ராஸ்
அதே கருத்தைத் தனியாகக் கண்டறிந்ததன் மூலம் 1872-இல்தான்
கணிதவியலாளர்கள் இதை முதன்முதலில் அறிந்தனர்.)\footnote{இதன் வரலாறு,
\href{http://en.wikipedia.org/wiki/Weierstrass_function}{வெயர்ஸ்ட்ராஸ்
சார்பு} பற்றிய விக்கிப்பீடியா கட்டுரையின் மிக விரிவான அடிக்குறிப்புகளில்
ஆவணப்படுத்தப்பட்டுள்ளது.} இதைக் கண்டு வியக்காமல் இருப்பது கடினம்.
$|x|$ போன்ற, எல்லா இடங்களிலும் தொடர்ச்சியாக இருந்தாலும் ஒரு
குறிப்பிட்ட புள்ளியில் வகையிடத்தக்கதாக இல்லாத சார்புகளைக் கண்டுபிடிப்பது
எளிது. ஆனால் எல்லா இடங்களிலும் தொடர்ச்சியாக இருந்து,
\emph{எந்த இடத்திலும்} வகையிடத்தக்கதாக இல்லாத சார்பு முற்றிலும்
வேறுபட்டது. அத்தகைய சார்பை எப்படி வரைவீர்கள் என்று சற்று
யோசித்துப் பாருங்கள். அது தொடர்ச்சியாக இருக்க, தாளிலிருந்து பேனாவை
எடுக்காமல் வரைய இயல வேண்டும்; அதேவேளை அது எந்த இடத்திலும்
வகையிடத்தக்கதாக இல்லாமல் இருக்க, பேனாவின் திசையை இடைவிடாமல்
திடீரென மாற்ற வேண்டியிருக்கும்.

% SEGMENT OLP-0636-S02
மேலும் இத்தகைய வழக்கத்திற்கு மாறான எடுத்துக்காட்டுகள் தோன்றின.
1874 ஜனவரி 5 அன்று காண்டோர், டெடெகிண்டிற்கு எழுதிய கடிதத்தில்
பின்வரும் சிக்கலை முன்வைத்தார்:
\begin{quote}
ஒரு பரப்பின் (அதன் எல்லையும் சேர்ந்த ஒரு சதுரம் என்று வைத்துக்கொள்வோம்)
புள்ளிகளுக்கும், ஒரு கோட்டின் (அதன் முனைகளும் சேர்ந்த ஒரு நேர்கோடு
என்று வைத்துக்கொள்வோம்) புள்ளிகளுக்கும் ஒன்றுக்கொன்று ஒருமையான
பொருத்தத்தை அமைக்க முடியுமா? அதாவது, பரப்பின் ஒவ்வொரு புள்ளிக்கும்
கோட்டில் ஒரு புள்ளி பொருந்தவும், மறுதிசையில் கோட்டின் ஒவ்வொரு
புள்ளிக்கும் பரப்பில் ஒரு புள்ளி பொருந்தவும் முடியுமா?

இப்போதும் இக்கேள்விக்குப் பதிலளிப்பது மிகவும் கடினம் என்றே எனக்குத்
தோன்றுகிறது. இருந்தாலும், இங்கும் \emph{முடியாது} என்று சொல்ல
மனமே தூண்டுகிறது; அதற்கான நிறுவல் ஏறக்குறைய தேவையற்றது என்றுகூடக்
கருதத் தோன்றுகிறது. [\citealt{Gouvea2011}-இல் மேற்கோள்]
\end{quote}
ஆனால் 1877-இல் தாம் தவறாக இருந்ததை காண்டோர் நிறுவினார்.
உண்மையில், ஒரு கோட்டிலும் ஒரு சதுரத்திலும் துல்லியமாக ஒரே எண்ணிக்கையிலான
புள்ளிகள் உள்ளன. 1877 ஜூன் 29 அன்று டெடெகிண்டிற்கு
“\emph{je le vois, mais je ne le crois pas}” என்று எழுதினார்;
அதாவது, “அதை நான் காண்கிறேன்; ஆனால் நம்ப முடியவில்லை.”
கணிதத்தின் பொதுவாகச் சொல்லப்படும் வரலாற்றில், இந்தப் புதிய முடிவுகள்
அக்காலக் கணிதவியலாளர்களுக்கு எவ்வளவு \emph{உண்மையாகவே நம்பமுடியாதவையாக}
இருந்தன என்பதைக் காட்டுவதாக இது அடிக்கடி கொள்ளப்படுகிறது.
(அந்தக் கடிதப் பரிமாற்றம் \citet{Gouvea2011}-இல் தரப்பட்டுள்ளது;
\olref[his][set][mythology]{sec} பகுதியில் மீண்டும் அதைப் பார்ப்போம்.
காண்டோரின் நிறுவலின் சுருக்கம் \olref[his][set][cantorplane]{sec}
பகுதியில் தரப்படுகிறது.)

% SEGMENT OLP-0636-S03
காண்டோரின் முடிவால் தூண்டப்பட்ட பியானோ, ஒரு கோட்டை ஒரு தளத்தின்மீது
\emph{இடைவெளியின்றி} படமாக்க முடியுமா என்று சிந்திக்கத் தொடங்கினார்.
அது \emph{வெளியை நிரப்பும் வளைவாக} இருக்கும். \citeyear{Peano1890}-இல்
பியானோ அத்தகைய வளைவை அமைத்தார். இது உண்மையிலேயே உள்ளுணர்வுக்கு
மாறானது: யூக்ளிட் ஒரு கோட்டை “அகலமற்ற நீளம்” (முதல் நூல்,
வரையறை 2) என்று வரையறுத்திருந்தார். ஆனால் ஒரு கோட்டைத் தக்கவாறு
மடக்குவதன் மூலம் அதன் நீளம் அகலத்தை நிரப்பலாம் என்று பியானோ
காட்டியிருந்தார். \citeyear{Hilbert1891}-இல் ஹில்பர்ட், சற்றே
எளிதாகப் புரியக்கூடிய மற்றொரு வெளிநிரப்பு வளைவையும் அதைக்
காட்டும் சில படங்களையும் விவரித்தார். அந்த வளைவு படிப்படியாக
அமைக்கப்படுகிறது; அதன் முதல் ஆறு நிலைகள் இவை:
\begin{center}
\begin{tikzpicture}
\begin{scope}[xshift=20pt, yshift=20pt]
	\draw [oldiagcolorC, l-system={Hilbert curve, step=40pt, angle=90, axiom=L, order=1}] lindenmayer system;
\end{scope}
\begin{scope}[xshift=110pt, yshift=10pt]
	\draw [oldiagcolorC, l-system={Hilbert curve, step=20pt, angle=90, axiom=L, order=2}] lindenmayer system;
\end{scope}
\begin{scope}[xshift=205pt, yshift=5pt]
	\draw [oldiagcolorC, l-system={Hilbert curve, step=10pt, angle=90, axiom=L, order=3}] lindenmayer system;
\end{scope}
\begin{scope}[xshift=2.5pt, yshift=-97.5pt]
	\draw [oldiagcolorC, l-system={Hilbert curve, step=5pt, angle=90, axiom=L, order=4}] lindenmayer system;
\end{scope}
\begin{scope}[xshift=101.25pt, yshift=-98.75pt]
	\draw [oldiagcolorC, l-system={Hilbert curve, step=2.5pt, angle=90, axiom=L, order=5}] lindenmayer system;
\end{scope}
\begin{scope}[xshift=200.625pt, yshift=-99.375pt]
	\draw [oldiagcolorC, l-system={Hilbert curve, step=1.25pt, angle=90, axiom=L, order=6}] lindenmayer system;
\end{scope}
\draw[gray] (0pt,0pt) rectangle (80pt, 80pt);
\draw[gray] (100pt,0pt) rectangle (180pt, 80pt);
\draw[gray] (200pt,0pt) rectangle (280pt, 80pt);
\draw[gray] (0pt,-100pt) rectangle (80pt, -20pt);
\draw[gray] (100pt,-100pt) rectangle (180pt, -20pt);
\draw[gray] (200pt,-100pt) rectangle (280pt, -20pt);
\end{tikzpicture}
\end{center}

% SEGMENT OLP-0636-S04
இப்போது துல்லியமான பொருள் பெற்றுவிட்ட எல்லை என்ற கருத்தின்படி,
முடிவிலான நிலையை நோக்கிச் சென்றால் முழுச் சதுரமும் அடர்த்தியான
!!{colorC} நிறத்தில் நிரம்பிவிடும். மேலும், ஹில்பர்ட்டின் வளைவு
எல்லா இடங்களிலும் தொடர்ச்சியானது; எந்த இடத்திலும்
வகையிடத்தக்கதல்ல. முடிவிலா எல்லையில் சார்பு ஒவ்வொரு கணமும்
திடீரெனத் திசை மாறுகிறது என்பது இதற்கான உள்ளுணர்வு விளக்கம்.
(ஹில்பர்ட்டின் கட்டுமானத்தை \olref[his][set][hilbertcurve]{sec}
பகுதியில் மேலும் விரிவாகக் காண்போம்.)

இத்தகைய வழக்கத்திற்கு மாறான வடிவியல் கட்டுமானங்கள், நல்லதோ
கெட்டதோ, வடிவியல் உள்ளுணர்வை நம்புவதைச் சந்தேகிக்க ஒரு
காரணமாகக் கொள்ளப்பட்டன. அவை கணக் கோட்பாட்டில் உச்சம் பெற்ற
புதிய கணிதச் செய்முறைக்கான \emph{பரப்புரையாக} மாறின.
பின்னர் இத்துறையைச் சுற்றி உருவாக்கப்பட்ட கதைகளில், இந்த முடிவுகள்
முற்றிலும் துல்லியமானவையும் அதே அளவு அதிர்ச்சியளிப்பவையும்
என்று மீண்டும் மீண்டும் கூறப்பட்டது. ஆகவே அவை இரு நோக்கங்களுக்குப்
பயன்பட்டன: வடிவியல் உள்ளுணர்வை மட்டும் நம்பக் கூடாது என்ற
எச்சரிக்கையாகவும், புதிய சிந்தனை முறைகளின் வளத்தை வெளிப்படுத்தும்
காட்சியாகவும்.

\end{document}
